\section{Integral annihilators and reflection for all four pure quartic labels}
\label{sec:quartic-abelian-reflection-strict}

\paragraph{Original labels and auxiliary field.}
Put $p=269$, $K=\mathbb Q(\sqrt{-3})$ and
$H=\mathbb Q(\zeta_{48p})$.
This fixed field is used only for residual restriction and class
groups; it does not replace any prescribed growing host.
Fix $i\in\mathbb Z_p$ with $i^2=-1$ and $\bar i=82$.
Let $\phi$ be the primitive odd quartic character modulo sixteen
with $\phi(13)=i$, $\phi(5)=-i$ and kernel $\{1,7\}$.
For $\sigma=1,-1$ and $a=0,1$ set
\[
              \delta_{\sigma,a}=\phi^\sigma\chi_{-4}^{\,a}|_{G_K}.
 \tag{G170.1}
\]
These are four distinct order-four characters, all already
available on the prescribed conductor-sixteen part of the hosts.
Write $c_p$ for the cyclotomic character and $\omega_p$ for
its Teichmuller part.  For a discrete coefficient $D$ use
\[
 \begin{gathered}
 S_{\rm str}(K,D)=\ker\bigl(H^1_{\rm cts}(G_K,D)
                                  \longrightarrow\mathcal L(D)\bigr),\\
 \mathcal L(D)=\prod_{v\mid p}H^1(G_{K_v},D)
                   \times\prod_{v\nmid p,\infty}H^1(I_v,D).
 \end{gathered}
 \tag{G170.2}
\]

\begin{theorem}[AC-G170: broad strict vanishing for the four pure quartic labels]
\label{AC-G170}
Define four exceptional sets and their complements by
\[
 \begin{aligned}
 \mathcal E_{1,0}&=\{243,248\},&
 \mathcal E_{-1,0}&=\{19,24\},\\
 \mathcal E_{1,1}&=\{11,60,229\},&
 \mathcal E_{-1,1}&=\{38,207,256\},\\
 \mathcal R_{\sigma,a}&=\{0,1,\ldots,267\}\setminus\mathcal E_{\sigma,a}.
 \end{aligned}
 \tag{G170.3}
\]
Let $E/\mathbb Q_p$ be any finite extension with its original
ring of integers $O$, allowing arbitrary absolute ramification.
Let $m\in\mathbb Z$, $e\geq1$, and let
$\kappa:G_K\to1+\mathfrak m_O$ be a continuous character
unramified at every finite place away from $p$.
If $r\equiv e+m\pmod{268}$ belongs to $\mathcal R_{\sigma,a}$, then
\[
 \begin{gathered}
 D=(E/O)(\delta_{\sigma,a}\omega_p^{-m}\kappa^{-1}c_p^{-e}),\\
                            S_{\rm str}(K,D)=0.
 \end{gathered}
 \tag{G170.4}
\]
The residual strict groups vanish over every finite field of
characteristic $p$ for the same residues.
In particular $r=0$ is included in both conclusions.
There are $266,266,265,265$ good residues for the four labels,
respectively, and $1062$ good label--residue pairs in total.
No nonvanishing assertion is made for the ten uncovered pairs.
For $m=0,\kappa=1$ the finite labels are precisely (G170.1),
with trivial cubic part.  Positive derived indices and the
separate actual-window comparison are not asserted zero.
\end{theorem}

\begin{proof}
\proofstep{1. The complete new character sector and primitive conductors.}
Put $G=\operatorname{Gal}(H/\mathbb Q)$ and
$J=\operatorname{Gal}(H/K)$.
Then
\[
 |G|=\varphi(48)\varphi(p)=4288,\qquad
 |J|=2144,\qquad p\nmid|G|.
 \tag{G170.5}
\]
The exponent of $G$ is $268=p-1$, so all its character
values lie in $\mathbb Z_p^\times$ and reduce without loss
of order to $\mathbb F_p^\times$.
Let $\omega$ be the Teichmuller lift of its action on $\mu_p$;
on $\sigma_n:\zeta_{48p}\mapsto\zeta_{48p}^{n}$ it reduces
to $n\pmod p$.

The eight tame characters with exact order four at two are
$\phi^\sigma\chi_{-4}^{\,a}\chi_{-3}^{\,b}$,
where $\sigma=\pm1$ and $a,b\in\{0,1\}$.
They are distinct: modulo sixteen the four characters
$\phi^\sigma\chi_{-4}^{\,a}$ are the four primitive quartics,
and the conductor-three factor is independent.
The other eight tame characters have order at most two and
are not part of the new finite table.
Consequently the complete quartic-at-two sector of $G$ is
\[
 \begin{gathered}
 \chi_{\sigma,a,b,k}
       =\phi^\sigma\chi_{-4}^{\,a}\chi_{-3}^{\,b}\omega^k,\\
 \sigma=\pm1,\quad a,b=0,1,\quad0\leq k<268.
 \end{gathered}
 \tag{G170.6}
\]
It contains $2144$ characters.
Their parity is $(-1)^{1+a+b+k}$; exactly $1072$ are odd,
namely those with $k+a+b$ even.
Their primitive conductors are
\[
 f_{a,b,k}=16\cdot3^b\quad(k=0),\qquad
 f_{a,b,k}=16\cdot3^b p\quad(k\ne0).
 \tag{G170.7}
\]
Indeed $\phi(9)=-1$ while $\chi_{-4}(9)=1$, so the
two-part is exactly sixteen even after the quadratic twist.
The coprime conductor-three factor is present exactly when
$b=1$, and a nontrivial character modulo the prime $p$ has
conductor $p$.  Thus every primitive conductor in (G170.7)
divides $48p$, including the four odd rows with $k=0$.

Put $A_H=\operatorname{Cl}(H)\otimes\mathbb Z_p$.
For a character $\chi$ write
\[
 e_\chi=|G|^{-1}\sum_{g\in G}\chi(g)^{-1}g,
                     \qquad A_H(\chi)=e_\chi A_H.
 \tag{G170.8}
\]
These are integral $p$-adic idempotents and give an exact
direct-sum decomposition, also after reduction or scalar extension.

\proofstep{2. Exact integral operators on all odd quartic rows.}
For $F=\mathbb Q(\zeta_f)$, let
\[
 \begin{aligned}
 \theta_f&=\frac1f\sum_{\substack{1\leq n<f\\(n,f)=1}}
                              n\sigma_n^{-1},\\
 S_{f,c}&=(c-\sigma_c)\theta_f\\
       &=\sum_{\substack{1\leq n<f\\(n,f)=1}}
                   \left\lfloor\frac{cn}{f}\right\rfloor\sigma_n^{-1}
                   \qquad((c,f)=1).
 \end{aligned}
 \tag{G170.9}
\]
The equality follows by reindexing multiplication by $c$:
the coefficient in $\sigma_c\theta_f$ is $\{cn/f\}$.
It proves integrality before reduction modulo $p$.
Sharifi's Definition~4.5.1, Lemma~4.5.2 and Theorem~4.5.6
state that the integral intersection of
$\mathbb Z[\operatorname{Gal}(F/\mathbb Q)]\theta_f$ with
the integral group ring annihilates $\operatorname{Cl}(F)$.
The source's generator $\sigma_c-c$ differs only by sign.
Thus every $S_{f,c}$ annihilates this ordinary class group.

For (G170.6) at its primitive conductor, the eigenvalue reduces to
\[
 \begin{aligned}
 s_{\sigma,a,b,k}(c)
 &=\sum_{\substack{1\leq n<f\\(n,f)=1}}
 \left\lfloor\frac{cn}{f}\right\rfloor
 \bar\phi(n)^{-\sigma}\\
 &\hspace{20mm}\cdot\chi_{-4}(n)^a\chi_{-3}(n)^b n^{-k}
                      \pmod p.
 \end{aligned}
 \tag{G170.10}
\]
For $k=0$ the last factor is one and the primitive modulus
does not include $p$.  For $k\ne0$, every summand is prime
to $p$ and the inverse power is defined.
Both the inverse quartic exponent and $-k$ come from
$\sigma_n^{-1}$; the quadratic factors are their own inverses.
A nonzero residue makes the exact $p$-adic eigenvalue a unit
and kills the corresponding $p$-primary class component of $F$.
No division by a $p$-divisible conductor in a residue field occurs.

The complete finite certificate evaluates (G170.10) for every
one of the $1072$ odd rows and every $2\leq c<32$ prime to
its primitive conductor, a total of $13400$ exact values.
Its character values are determined by
\[
 \begin{array}{c|rrrrrrrr}
 n\bmod16&1&3&5&7&9&11&13&15\\ \hline
 \phi(n)&1&-i&-i&1&-1&i&i&-1,
 \end{array}
 \tag{G170.11}
\]
with $i\equiv82$, the ordinary residue rules for
$\chi_{-4},\chi_{-3}$, and $n^{-k}$ modulo $269$.
For $k\ne0$ the verifier first groups the floor coefficients
by $n\pmod p$, orders the $268$ nonzero residues as $2^j$,
and evaluates the resulting polynomial at $2^{-k}$ by Horner's rule.
The integer $2$ is checked to be a primitive root.
The $k=0$ rows are computed separately at the smaller conductor.

Exactly $1067$ rows have a displayed nonzero witness.
The only odd rows without one in this finite list are
\[
 \begin{aligned}
 \mathcal Z^-=\{&(1,1,0,39),\ (1,1,1,208),\ (-1,0,0,244),\\
               &(-1,0,1,249),\ (-1,1,1,12)\}.
 \end{aligned}
 \tag{G170.12}
\]
The four coordinates are $(\sigma,a,b,k)$.
These zero rows are not a class-nonvanishing assertion and
do not assert that all integral Stickelberger operators vanish.

The primitive-field result passes to $H$ without extra Euler
factors.  Namely, for $U=\operatorname{Gal}(H/F)$,
\[
 N_{H/F}j_{H/F}=|U|,\qquad
             j_{H/F}N_{H/F}=\sum_{u\in U}u.
 \tag{G170.13}
\]
Since $|U|$ is prime to $p$, these identities identify
$A_F$ with $A_H^U$, equivariantly for the remaining Galois action.
For surjectivity, $x\in A_H^U$ satisfies
$j_{H/F}N_{H/F}x=|U|x$, and $|U|$ is invertible on the
finite $p$-group.  The character under consideration is
trivial on $U$, so its component is exactly the primitive
field component.  This is also the precise statement and
proof of Sharifi Lemma~1.3.5.
We have proved
\[
 A_H(\chi_{\sigma,a,b,k})=0
 \quad\text{for every odd quartic row outside }\mathcal Z^-.
 \tag{G170.14}
\]

\proofstep{3. Reflection, with the exact quartic inverse and local correction.}
We give the needed specialized reflection argument on this
actual $H$, rather than replacing it by a different field.
Sharifi Theorem~4.4.2 gives, for an odd character $\chi$,
\[
 r_p A_H(\omega\chi^{-1})-\delta_\chi
                                  \leq r_p A_H(\chi).
 \tag{G170.15}
\]
Its correction can be nonzero only for $\chi=\omega$ and
only when $H(\mu(H)^{1/p})/H$ is unramified.
The base field is the totally real field $\mathbb Q$, the
extension $H/\mathbb Q$ is abelian of degree prime to $p$,
and it contains $\mu_p$ and every character field in use.
Thus the exact source hypotheses hold.

For completeness, let $B\subset H^\times$ consist of radicals
whose $p$th roots generate unramified extensions of $H$.
Kummer theory and the ordinary unramified class-field pairing
give the $G$-equivariant exact sequence
\[
 0\longrightarrow(B\cap\mathcal O_H^\times)/\mathcal O_H^{\times p}
 \longrightarrow\operatorname{Hom}(A_H/pA_H,\mu_p)
 \longrightarrow A_H[p].
 \tag{G170.16}
\]
There is no surjectivity assertion at the final term.
An unramified radical has all valuations divisible by $p$;
writing $(x)=\mathfrak a^p$ sends it to $[\mathfrak a]$.
Its kernel consists exactly of unit radicals modulo $p$th powers.
This is the sequence proved in Sharifi Lemma~1.4.6 using
Proposition~1.4.4, with the ordinary class group, not a group
with primes above $p$ inverted.

Because $H$ is CM, the minus part of its units modulo torsion
is zero.  Its only possible odd unit line modulo $p$ is
$\mu(H)/\mu(H)^p$, of character $\omega$; this is the
equivariant Dirichlet argument of Sharifi Lemma~1.4.12.
The $p$-primary roots of unity in $H$ are exactly $\mu_p$,
since its degree is prime to $p$.
At $p$, the conductor-forty-eight part is unramified of degree
four: $269\equiv29\pmod{48}$, $29^2\equiv25$ and
$29^4\equiv1$, with smaller positive powers nontrivial.
Each completion of $H$ therefore has indices
\[
                   e=p-1=268,\qquad f=4.
 \tag{G170.17}
\]
It is the compositum of $\mathbb Q_p(\mu_p)$ and that
unramified field.  There are four primes above $p$.
Adjoining $\mu_{p^2}$ is totally ramified of degree $p$ over
each completion: the cyclotomic Eisenstein polynomials after
$X=1+T$ give ramification degrees $p-1$ and $p(p-1)$,
and unramified base change preserves their ratio.
The prime-to-$p$ roots of unity already have $p$th roots in $H$.
Hence every nonzero class of $\mu(H)/\mu(H)^p$ gives a
ramified Kummer extension, and the odd kernel in (G170.16)
is zero, including its $\omega$ component.
In particular $\delta_\chi=0$ in (G170.15).

For an even character $\nu$, the dual of its class component
in the middle term of (G170.16) has character
$\omega\nu^{-1}$, since
\[
                        (g f)(x)=g\,f(g^{-1}x).
 \tag{G170.18}
\]
The zero odd kernel gives an injection into
$A_H[p](\omega\nu^{-1})$ and therefore
\[
 A_H(\omega\nu^{-1})=0\quad\Longrightarrow\quad A_H(\nu)=0.
 \tag{G170.19}
\]
A finite $p$-group with zero quotient modulo $p$ is zero;
no identification of $A_H[p]$ with $A_H/pA_H$ is assumed.
For the characters (G170.6), the reflection map is exactly
\[
          R(\sigma,a,b,k)=(-\sigma,a,b,1-k\bmod268).
 \tag{G170.20}
\]
The inversion of $\sigma$ is essential for quartic characters.
Combining (G170.14) and (G170.19), every quartic-at-two
class component vanishes outside the ten possible rows
$\mathcal Z^-\cup R(\mathcal Z^-)$.
This statement is restricted to that sector; it makes no
additional assertion about the quadratic-at-two components.

\proofstep{4. Strict residual classes select the original class character.}
Let $k$ be any finite field of characteristic $p$, and let
$\alpha=\bar\delta_{\sigma,a}\bar\omega_p^{-r}$ on $G_K$.
The field $H$ splits $\alpha$ and $[H:K]=2144$ is prime to $p$.
Restriction injects $H^1(G_K,k(\alpha))$ into
$H^1(G_H,k)$ with the twisted quotient action.
If $z$ is a cocycle, its restriction is a homomorphism and
the cocycle identity gives
\[
           z(\tau g\tau^{-1})=\alpha(\tau)z(g)
                    \quad(g\in G_H,\ \tau\in G_K).
 \tag{G170.21}
\]
If the original class is strict, this homomorphism vanishes
on every inertia group of $G_H$ and on every full decomposition
group above $p$.  The local coboundary is literally zero on
these subgroups because the coefficient is trivial there.
The homomorphism therefore factors through the class group
with primes above $p$ removed, which is a quotient of the
ordinary class group.  There is consequently an injection
\[
 S_{\rm str}(K,k(\alpha))\lhook\joinrel\longrightarrow
       \operatorname{Hom}_{\mathbb F_p[J]}
                         (A_H/pA_H,k(\alpha)).
 \tag{G170.22}
\]
Only this injection is asserted.  Galois equivariance of the
Artin map, in Sharifi Proposition~1.3.2, makes (G170.21)
select the class character $\alpha$, not its inverse.
This is distinct from the duality in (G170.18).

The two rational extensions of this character of $J$ are
\[
               \chi_{\sigma,a,0,-r},\qquad
               \chi_{\sigma,a,1,-r},
 \tag{G170.23}
\]
with exponents read modulo $268$.
They differ by $\chi_{-3}$ and have opposite parity.
Integral idempotents show that only these two components
can contribute to (G170.22), also after extension to $k$.
For each $b=0,1$, the exact odd source row needed is
\[
 \begin{cases}
 (\sigma,a,b,-r),&-r+a+b\text{ even},\\
 (-\sigma,a,b,1+r),&-r+a+b\text{ odd}.
 \end{cases}
 \tag{G170.24}
\]
The first line uses an odd component directly; the second
uses (G170.19), not an unchanged-sign quadratic formula.

\proofstep{5. Every residue and all ten uncovered pairs.}
The exact translation verifier considers all $268$ residues
for each of the four original labels and both rows in
(G170.24).  The five possible odd failures give exactly
\[
 \begin{array}{c|c|c}
 (\sigma,a,b,k)&\text{direct }(\sigma,a;r)&\text{reflected }(\sigma,a;r)\\ \hline
 (1,1,0,39)&(1,1;229)&(-1,1;38)\\
 (1,1,1,208)&(1,1;60)&(-1,1;207)\\
 (-1,0,0,244)&(-1,0;24)&(1,0;243)\\
 (-1,0,1,249)&(-1,0;19)&(1,0;248)\\
 (-1,1,1,12)&(-1,1;256)&(1,1;11).
 \end{array}
 \tag{G170.25}
\]
The notation $(\sigma,a;r)$ denotes an original label and
its effective residue.  Direct use has $r=-k$, while reflection
has original sign $-\sigma$ and $r=k-1$.
All ten pairs in the last two columns are distinct.
For every other pair, both required odd source rows have
nonzero integral operator witnesses.  Thus both class
components in (G170.23) vanish, and (G170.22) proves
residual strict zero for precisely the sufficient sets
$\mathcal R_{\sigma,a}$ in (G170.3).
In particular each of the four $r=0$ pairs is supported.
The verifier records every source row and witness, rather
than inferring the residue assertion merely from its cardinality.

\proofstep{6. The original coefficient lift, including residue zero.}
Let $\varpi$ be a uniformizer of the original $O$ and let $D$
be the module in (G170.4).
It is divisible, with $D[\varpi]\simeq k(\alpha)$.
If $r\ne0$, the finite conductor-sixteen character is
unramified at $p$, and the residual inertia character is
$\bar c_p^{-r}\ne1$ since $K_p/\mathbb Q_p$ is unramified.
An inertia scalar minus one is an $O$-unit, so
$D^{G_{K_p}}=0$.

For $r=0$ one must instead use the actual unramified finite
character.  At the rational prime $269$,
\[
 \phi(269)=\phi(13)=i,
 \qquad \chi_{-4}(269)=1,
 \qquad \delta_{\sigma,a}(\operatorname{Frob}_{K_p})
                         =(i^\sigma)^2=-1.
 \tag{G170.26}
\]
The last identity uses the unramified quadratic local extension
$K_p/\mathbb Q_p$.  It is unaffected by choosing arithmetic
or geometric Frobenius, since $(-1)^{-1}=-1$.
At residue zero the cyclotomic and Teichmuller product reduces
to one on the full local group, and $\kappa$ also reduces to one.
Thus a Frobenius scalar minus one reduces to $-2\ne0$,
again proving $D^{G_{K_p}}=0$.
This is the additional original-coefficient argument at $r=0$;
it is not inferred from residual vanishing alone.

Away from $p$, the cyclotomic, Teichmuller and $\kappa$
factors are unramified.  The remaining finite inertia action
has order dividing four, prime to $p$, so its invariants are
either zero or the whole divisible module $D$.
For every local group $L$ in (G170.2), therefore,
\[
                            D^L/\varpi D^L=0.
 \tag{G170.27}
\]
The multiplication-by-$\varpi$ long exact sequence gives
local injectivity $H^1(L,D[\varpi])\to H^1(L,D)$.
A strict global class killed by $\varpi$ lifts to a residual
class, which remains strict by local injectivity and is zero
by Step~5.
Every continuous cocycle from compact $G_K$ into discrete
$E/O$ has finite image killed by a power of $\varpi$.
A nonzero strict class would then yield a nonzero one killed
by $\varpi$, a contradiction.  This proves (G170.4) on the
original, arbitrarily ramified coefficient ring.
\end{proof}

\paragraph{Actual labels and exact finite-weight counts.}
For $m=0,\kappa=1$ the four labels (G170.1) factor through
$\mathbb Q(\zeta_{16})$, which is contained in every host
of (FG.3) with $h\geq2$.
They remain distinct over $K$, since
$K\cap\mathbb Q(\zeta_{16})=\mathbb Q$.
The cases $a=0$ are the two odd quartics, and $a=1$ gives
the two even quartics; none is quadratic or trivial.
All have trivial cubic part.  Neither $\omega_p$ nor the
general lift $\kappa$ is added as an actual finite host label.

For positive integer $T$, the exact number of weights
$1\leq e<T$ supplied for label $(\sigma,a)$ is
\[
 (T-1)-\sum_{r\in\mathcal E_{\sigma,a}}
                   \left\lfloor\frac{T+267-r}{268}\right\rfloor.
 \tag{G170.28}
\]
Every exceptional residue is nonzero, so the summands count
the missing progressions without an $e=0$ contribution.
This finite count does not remove any derived-index comparison
cost.  The actual all-prime remainder and common-window
transport remain separate, as do mixed realization and H04.4.

\paragraph{Source and finite-evidence boundaries.}
The finite files are \path{compute_quartic_odd_annihilators.py},
\path{quartic_odd_annihilator_table.json},
\path{verify_quartic_residue_translation.py}, and
\path{quartic_residue_translation.json}.
They cover only the new quartic-at-two sector, not the
quadratic table of AC-G168.
The complete applied Sharifi Theorem~4.4.2 proof and the
Chapter~1 norm/extension, unramified Kummer, CM-unit and
Galois-action passages were freshly read in the retained
primary text.  The complete Definition~4.5.1 statement,
Lemma~4.5.2 proof and Theorem~4.5.6 proof retain the exact
inverse group-ring convention already checked for AC-G166.
The specialized reflection and strict-character arguments
are written out above, with the changed local degree, quartic
inverse and residue-zero lift all explicit.
The full AC-G168 proof was read as a comparison, without
treating its different-field theorem as a black-box substitute.
Deep class-field theory and the previously audited SGA~1,
Expos\'e~V, \S\S4 and 7 and Stacks tag~03PK dictionaries remain
at their stated source scopes.  No new higher regulator,
unlisted coefficient-prime bound, motive or Lean execution
is asserted.
